Detecting trace impurities in crowded and concentrated fluid environments requires balancing two competing effects: increasing number of searching probe particles accelerates initial target discovery, yet sufficiently high probe density eventually hinders search through steric exclusion and caging effects. Prior first-passage search theory has largely treated the number of independent searchers as a quantity to be maximized without limit, obscuring the density at which additional probes stop helping and start hindering detection. This gap is addressed via an in silico Langevin dynamics study, implemented using the LAMMPS particle-based simulation package, of a two-dimensional, two-species Brownian mixture of \SI{2.0}{\micro\meter} polystyrene probe particles searching for smaller, \SI{1.0}{\micro\meter} foreign target particles under Weeks-Chandler-Andersen excluded-volume repulsion and periodic boundary conditions. Across six system sizes ($N_\text{total} = 10$ to $200$), the mean first-passage detection time $\langle T \rangle$ was mapped against probe number $N_p$ and fit to a two-term ``Caldera" model balancing diffusion-limited decay against crowding-limited growth, while the ergodic product $N_p \times \langle T \rangle$ was tracked as an independent diagnostic of the breakdown of ideal, non-interacting searcher scaling. The \textit{caldera} minimum consistently locates an optimal probe density near an area fraction of $\phi^*\approx29-31\%$ for five of the six system sizes studied ($N_\text{total} = 30$--200), with a modestly elevated optimum ($\phi^* \approx 34\%$) at the smallest, most coarsely sampled system ($N_\text{total} = 10$); in every case the optimum sits well below the geometric close-packing ceiling, and the ergodic product flags this crowding onset tens of probes before it becomes visible in $\langle T \rangle$ alone. The resulting caldera profiles were further classified into a four-part typology --- Incipient, Shield, and Steep-Walled calderas reflecting how sharply detection time rises once the crowding-limited regime is entered, and a fourth, qualitatively distinct bifurcated caldera shape. These results provide concrete, system-size-independent probe-density guidance for point-of-care microfluidic and colloidal impurity-detection analysis.

% to edit - will add Doctor Saloma suggestionsmes